\section{Result, scope, and relation to the repository}\label{sec:results}

\subsection{The concrete problem addressed}

Let $\mathcal T$ be a finite triangulation of a compact connected orientable
three-manifold whose boundary is one torus. Let $x\in\Z_{\geq0}^{7t}$ be an
admissible normal vector, where $t$ is the number of tetrahedra. Write $S(x)$
for the corresponding embedded normal surface. The input coordinates are
binary integers; neither $S(x)$ nor its components are supplied as expanded
cell complexes. Define
\[
  \D(x)=\#\{C\in\pi_0(S(x)):C\text{ is a disc with essential boundary}\}.
\]
All multiplicities in this definition count actual components. In particular,
$10^{100}$ parallel meridian discs contribute $10^{100}$, not one.

The previous native backend obtains the number of components, orientability
counts, boundary-curve count, and total Euler characteristic through interval
orbits, with optional additional classification of closed and boundary-touching
components. Its positive disc predicate requires the \emph{whole supplied
surface} to be connected. The maintained boundary-classification discussion
explicitly leaves componentwise disc identification unresolved
\cite{ProveItSnapshot}. The new operation computes $\D(x)$ for a disconnected
surface and supplies replayable evidence for the answer.

This is a sharply delimited algorithmic improvement. A positive answer proves
that the supplied manifold has compressible boundary. If that manifold has
independently been identified with the exterior of an input knot in $S^3$, it
certifies unknottedness. A zero answer only eliminates this particular vector.
It does not eliminate other normal surfaces, prove boundary incompressibility,
or establish knottedness. The delivered code does not silently add any of
these missing implications.

\subsection{Main statements}

\begin{theorem}[Three-weight disc query]\label{thm:main-disc}
For the validated input class above, there is an algorithm that computes
$\D(x)$ and a compressed histogram of componentwise Euler characteristic and
boundary homology using three integer weights per orbit representative.
Its running time, output length, and certificate length are polynomial in
$t$ and the maximum input-coordinate bit length. Its certificate can be
checked without invoking its orbit-discovery scheduler or its weight-transport
implementation.
\end{theorem}

The topology in this theorem is proved in \cref{sec:geometry}; the compressed
algorithm and its bit accounting are established in \cref{sec:weighted-orbits}.
The checker shares the finite-manifold input validator and the deterministic
cell-weight construction. Independence therefore concerns discovery and
arithmetic transport, not two unrelated implementations of the entire
topological input model.
A further refinement to two weights after certified ambient homology
preprocessing is proved in \cref{sec:two-weight-refinement}. That variant is
not implemented or included in the reported timings.

\begin{theorem}[Primitive quadrilateral core]\label{thm:main-core}
For each global vertex $v$, let $\ell_v$ be the least triangle coordinate
among the tetrahedral corners at $v$, and let $L_v$ be its vertex-link vector.
If $x$ has a nonzero quadrilateral coordinate, let $g$ be the gcd of its
quadrilateral coordinates and set
\[
 P=\frac{x-\sum_v\ell_v L_v}{g},\qquad
 Q=\max_j\frac{x_{q_j}}{g}.
\]
Then $P$ is an admissible nonnegative integral normal vector and
\[
                  \D(x)=g\D(P).
\]
Every coordinate of $P$ is at most $\max(1,4t-1)Q$. If all quadrilateral
coordinates vanish, $\D(x)=0$ and no orbit computation is needed.
\end{theorem}

The full proof, including the effect of one-sided components under scaling,
appears in \cref{sec:quadrilateral-core}. The central computational consequence is that
the expensive orbit query can use
\[
                    O(\log(t+1)+\log(Q+1))
\]
bits per geometric coordinate. Large common multiplicity and large independent
vertex-link offsets still have to be read and processed, but they disappear
from the subsequent orbit search. This remains useful when the gcd of the
\emph{original seven-coordinate vector} is one.

\subsection{Which parts are established antecedents?}

The weighted orbit method and polynomial component recovery are classical
Agol--Hass--Thurston results. Canonical removal of vertex links and conversion
between quadrilateral and standard coordinates are classical Q-theory.
We use them openly rather than describe them as new topology
\cite{AHT2006,Tollefson1998,Burton2009}. The contribution here is the implemented
specialization, its source-bound certificate, its exact disc-count reduction,
and its measured improvement against more expensive routes to the same query.
We do not claim that the specialized mathematical statements have no earlier
appearance in the literature.

\begin{center}
\begin{tabular}{>{\raggedright\arraybackslash}p{0.30\linewidth}>{\raggedright\arraybackslash}p{0.59\linewidth}}
\toprule
Ingredient & Status in this continuation \\
\midrule
Weighted interval orbits & Classical algorithm; newly implemented as a
checked extension of the maintained repository trace. \\
Vertex-link canonicalization & Classical construction, used with explicit
binary-cost and quadrilateral-content accounting. \\
Three statistics for disc detection & A proved specialization that avoids
recovering $7t$ weights or counting orientation and boundary covers. \\
Two weights with ambient homology & Proved refinement; implementation and
performance comparison remain future work. \\
Disconnected disc census & New native repository capability with finite
oracles and source-bound replay. \\
General quasi-polynomial recognition & Remains unimplemented and unproved
for the delivered pipeline. \\
\bottomrule
\end{tabular}
\end{center}

\subsection{Why aggregate invariants are insufficient}

There are two different pitfalls. First, even parallel meridian discs have
zero \emph{total} mod-two boundary class. Thus requiring a nonzero total
class rejects valid positive examples. Second, a surface can have total
Euler characteristic one and a nonzero total boundary class without containing
a disc. For example, a small interior vertex-link sphere disjoint from an
orientable genus-one spanning surface has total Euler characteristic
$2+(-1)=1$. The sphere has no boundary and the genus-one component is not a
disc. The corresponding finite normal-surface example is included in the
external audit.

The new query associates Euler characteristic and boundary class to the
\emph{same connected component}. Its correctness does not follow from a
test of the sums of these quantities. This distinction is the main reason
weighted orbits are needed.

\subsection{Exact API contracts}

The new APIs are additive modules under \code{fastunknot}; existing recognizer
defaults, certificate versions, and old aggregate APIs are unchanged. The
following table specifies the information computed by each route.

\begin{center}
\begin{tabular}{>{\raggedright\arraybackslash}p{0.29\linewidth}>{\raggedright\arraybackslash}p{0.12\linewidth}>{\raggedright\arraybackslash}p{0.47\linewidth}}
\toprule
Route & Dimension & Certified output \\
\midrule
\code{mode='disk'} & $3$ & Component histogram of $(\chi,h_1,h_2\bmod2)$,
total components, and essential-disc count. \\
\code{mode='summary'} & $5$ & The preceding data, normal-disc count, and
boundary intersection weight for each histogram type. \\
\code{mode='coordinates'} & $7t$ & Actual connected-component normal vectors,
grouped with binary multiplicities, and their disc predicates. \\
Reduced count API & $3$ on $P$ & Essential-disc count of the original $x$,
with a certified canonical reduction. \\
\bottomrule
\end{tabular}
\end{center}

Boundary intersection weight is the number of normal boundary edge points,
equal to the number of boundary normal arcs. It is \emph{not} the number of
boundary circles. The reduced count API intentionally makes no claim about
the total number of original components: that quantity is not homogeneous
in the presence of one-sided components. The full-coordinate route can
extract a concrete disc vector for a later cutting operation, but finding
and certifying that cut is outside this package.

\subsection{The quasi-polynomial target in the inspected sources}

The explicit March 2021 slides announce
$2^{O((\log n)^3)}=n^{O((\log n)^2)}$ time; the institutional announcement
instead uses $n^{c\log n}$. Both are quasi-polynomial, but the exponents are
different \cite{LackenbySlides,Oxford2021}. The July 2026 preprint provides a
hierarchy algorithm and polynomial verification results, while separating
iteration counts from step costs that have not all been specified sufficiently
for runtime analysis \cite{Lackenby2026}. The repository makes the same
distinction. Consequently this article treats quasi-polynomial time as the
global implementation target and states the local proved bounds explicitly.

The selected source snapshot is commit \baseline, whose full identity is
recorded in the provenance manifest. This pins the comparison despite active
development elsewhere in the repository. No remote repository changes are
required to inspect, reproduce, or integrate the accompanying archive.
